\chapter{Medical Tests}

\medterm{Abdominal Aortic Aneurysm Screening}-- Ultrasound screening for AAA recommended for men aged 65-75 with smoking history (USPSTF). A one-time abdominal ultrasound detects infrarenal aortic diameter \(\ge\)3 cm. Surveillance: 3.0-3.9 cm q3 years, 4.0-4.4 cm q12 months, 4.5-5.4 cm q6 months. Surgical referral at \(\ge\)5.5 cm or rapid enlargement (>1 cm/yr). Screening reduces AAA-related mortality by 50\%.

\medterm{Abdominal X-Ray}-- Plain radiograph of the abdomen (KUB: kidneys, ureters, bladder). Indications: suspected bowel obstruction (air-fluid levels, dilated loops), perforation (free air under diaphragm), constipation, foreign body, and renal calculi (90\% of stones are radiopaque). Sensitivity low for most intra-abdominal pathology. Uses: limited to specific clinical scenarios; CT is preferred for most acute abdominal conditions.

\medterm{ABO Blood Typing}-- Determination of a patient's ABO and Rh blood group prior to transfusion or organ transplantation. ABO system: A (A antigen, anti-B antibody), B (B antigen, anti-A antibody), AB (both antigens, no antibodies), O (no antigens, both antibodies). Rh: D antigen positive (85\%) or negative (15\%). Universal donor: O negative. Universal recipient: AB positive. Crossmatching confirms compatibility between donor and recipient.

\medterm{Acid-Fast Bacillus (AFB) Smear}-- Microscopic examination of sputum for acid-fast bacilli (Mycobacterium tuberculosis), using Ziehl-Neelsen or auramine-rhodamine stain. Positive smear indicates active pulmonary TB and correlates with infectiousness. Sensitivity: 50-80\% in culture-positive TB. Three specimens collected on consecutive mornings. Grading: negative, 1+ (1-9 per 100 fields), 2+ (1-9 per 10 fields), 3+ (1-9 per field), 4+ (>9 per field). Culture and NAAT are more sensitive.

\medterm{AFP (Alpha-Fetoprotein)}-- A tumor marker and prenatal screening analyte. Tumor marker: hepatocellular carcinoma (elevated in 60-80\%, >400 ng/mL highly specific), yolk sac tumors (germ cell tumors). Non-specific: cirrhosis, hepatitis, pregnancy. Prenatal: elevated in neural tube defects (maternal serum AFP). Screening for HCC: AFP + ultrasound q6 months in high-risk patients (cirrhosis, chronic HBV/HCV). AFP-L3 (isoform) improves specificity.

\medterm{Allergy Testing}-- Identification of IgE-mediated allergic sensitization. Skin prick test (SPT): allergen extract applied to skin via lancet, read after 15 min (wheal >3mm = positive). Intradermal test: more sensitive, used when SPT negative but suspicion remains. Specific IgE (formerly RAST): blood test measuring allergen-specific IgE levels; unaffected by antihistamines. Total IgE: elevated in atopic disease and parasitic infection. Indications: allergic rhinitis, asthma, food allergy, drug allergy, and anaphylaxis evaluation.

\medterm{Alvarado Score}-- Clinical scoring system for acute appendicitis. Components (MANTRELS): Migration of pain (1), Anorexia (1), Nausea/vomiting (1), Tenderness in RLQ (2), Rebound pain (1), Elevated temperature (1), Leukocytosis (2), Left shift (1). Scores: 1-4 (low probability, observe), 5-6 (moderate probability, CT or ultrasound), 7-10 (high probability, surgical consult). MASS (Modified Alvarado) removes left shift and uses 0-9.

\medterm{Ammonia}-- Measurement of serum ammonia (NH3). Normal: 15-45 mcg/dL (9-33 micromol/L). Elevated in hepatic encephalopathy (monitoring and diagnosis), urea cycle disorders (very high, >200 in neonates), Reye syndrome, and valproate therapy. Timing: must be collected in pre-chilled tube on ice and processed within 20 minutes (falsely elevated by delayed processing, smoking, exercise, and hemolysis). Used in initial evaluation of altered mental status with liver disease and suspected metabolic disorders.

\medterm{Amniocentesis}-- Sampling of amniotic fluid for prenatal diagnosis. Typically performed at 15-20 weeks gestation. Indications: advanced maternal age (\(\ge\)35), abnormal quad screen, family history of genetic disorders, and fetal lung maturity assessment. Analyzed for karyotype, FISH (aneuploidy of 13, 18, 21, X, Y), microarray, and AFP (neural tube defects). Procedure: ultrasound-guided needle aspiration of 15-20 mL of fluid. Risks: miscarriage (0.1-0.3\%), infection, and fetal injury.

\medterm{Amylase and Lipase}-- Pancreatic enzymes, elevated in acute pancreatitis. Lipase is more specific (pancreas-only) and remains elevated longer than amylase. Normal: amylase 30-110 U/L, lipase 10-140 U/L. Diagnosis: elevation >3x upper limit of normal, with compatible clinical presentation (epigastric pain, nausea/vomiting). Amylase also elevated in salivary gland disease, intestinal obstruction, perforated ulcer, and ectopic pregnancy. Lipase is the preferred single test for pancreatitis.

\medterm{Analysis of Urine}-- See Urinalysis.

\medterm{Angiography}-- Invasive imaging of blood vessels by injecting contrast material under fluoroscopy. Coronary angiography (cardiac catheterization): assesses coronary artery stenosis, guides PCI. Cerebral angiography: evaluates aneurysms, AVMs, and vasculitis. Peripheral angiography: evaluates PAD and guides intervention. Requires arterial access (femoral, radial). Complications: bleeding, pseudoaneurysm, contrast nephropathy, atheroembolism, and dissection. CTA and MRA have largely replaced diagnostic angiography.

\medterm{Angioplasty}-- Percutaneous coronary intervention (PCI) to open stenotic coronary arteries. Balloon angioplasty: inflation of balloon at the lesion to compress plaque. Bare metal stent (BMS): metal scaffold, restenosis 20-30\%. Drug-eluting stent (DES): elutes antiproliferative (sirolimus, paclitaxel, everolimus), restenosis 5-10\%. Dual antiplatelet therapy (DAPT: aspirin + P2Y12 inhibitor) required after stent placement. Indications: STEMI, NSTEMI, unstable angina, and stable angina refractory to medical therapy.

\medterm{Ankle-Brachial Index}-- Ratio of systolic BP at the ankle (dorsalis pedis or posterior tibial artery) to the brachial artery, measured with Doppler. Normal: 1.0-1.4. PAD: ABI \(\le\)0.90 (mild 0.7-0.9, moderate 0.4-0.7, severe <0.4). Non-compressible vessels (ABI >1.4, shows calcified vessels in diabetes). Uses: screening for PAD, wound healing assessment, and cardiovascular risk stratification. Toe-brachial index (TBI) if ABI unreliable.

\medterm{Antibiogram}-- A periodic summary of antimicrobial susceptibility patterns of bacterial isolates from a hospital or laboratory. Guides empiric antibiotic therapy by showing local resistance patterns (e.g., \% of E. coli susceptible to TMP-SMX). Updates annually. Clinicians use antibiograms to select appropriate empiric antibiotics while awaiting culture results. Cumulative antibiograms report susceptibilities for the most common pathogens.

\medterm{Anti-CCP Antibody}-- Anti-cyclic citrullinated peptide antibody, a highly specific serologic marker for rheumatoid arthritis (specificity >95\%, sensitivity 70\%). Positive anti-CCP is associated with more aggressive erosive RA. Appears early in disease and may precede clinical symptoms. Used in the ACR/EULAR classification criteria with RF and acute phase reactants. Testing recommended in any patient with unexplained inflammatory polyarthritis.

\medterm{Anti-dsDNA Antibody}-- Anti-double-stranded DNA antibody, highly specific for systemic lupus erythematosus (specificity 95\%, sensitivity 60-70\%). Elevated titers correlate with active disease, particularly lupus nephritis and systemic vasculitis. Serial monitoring used to assess disease activity and treatment response. In the ACR classification criteria for SLE. Other autoantibodies in SLE: anti-Smith (specific but low sensitivity), anti-Ro/SSA, anti-La/SSB, anti-RNP.

\medterm{Anti-Mullerian Hormone}-- AMH is a glycoprotein produced by granulosa cells of ovarian follicles, reflecting the ovarian reserve (quantity of remaining oocytes). Uses: assessment of ovarian reserve in infertility evaluation, predicting response to ovarian stimulation in IVF, diagnosis of PCOS (elevated), and monitoring for granulosa cell tumors. AMH declines with age and is undetectable after menopause. Unlike FSH, AMH is cycle-independent and reliable any day of the menstrual cycle.

\medterm{Antinuclear Antibody (ANA)}-- A screening test for autoimmune connective tissue diseases. Positive in SLE (sensitivity 95\%), Sjogren syndrome, systemic sclerosis, mixed connective tissue disease, and many other autoimmune conditions. Titer: 1:40 (low positive, common in healthy), 1:160 (clinically significant), 1:320 (strongly positive). Pattern: homogeneous (SLE, drug-induced), speckled (SLE, Sjogren, MCTD), nucleolar (scleroderma), centromere (limited scleroderma). Positive ANA requires further testing for specific autoantibodies.

\medterm{Arterial Blood Gas (ABG)}-- Direct measurement of arterial blood for pH, PaCO2, PaO2, HCO3, base excess, and oxygen saturation. Indications: respiratory failure assessment, acid-base disorder evaluation, ventilator management, and critical care. Normal values: pH 7.35-7.45, PaCO2 35-45 mmHg, PaO2 80-100 mmHg, HCO3 22-26 mEq/L, SaO2 95-100\%. Sample: radial artery (Allen test for collateral flow). Complications: hematoma, arterial spasm, thrombosis, and infection.

\medterm{Arthrocentesis}-- Aspiration of joint fluid for diagnostic or therapeutic purposes. Indications: septic arthritis (urgent, crystals), inflammatory arthritis (gout, pseudogout), hemorrhagic effusion (trauma, coagulopathy), and therapeutic drainage for symptomatic relief. Synovial fluid analysis: cell count/differential, Gram stain/culture, crystal analysis (polarized light microscopy: monosodium urate = needle-shaped, negatively birefringent; CPPD = rhomboid, weakly positively birefringent), and glucose/protein.

\medterm{Audiometry}-- Objective measurement of hearing thresholds across frequencies. Pure tone audiometry: air conduction (headphones) and bone conduction (mastoid oscillator) thresholds at 250-8000 Hz. Conductive hearing loss: abnormal air, normal bone (air-bone gap). Sensorineural: abnormal both. Speech audiometry: speech reception threshold (SRT) and word recognition score (WRS). Tympanometry: middle ear function (type A normal, B effusion, C negative pressure). Used for hearing loss diagnosis and hearing aid candidacy.

\medterm{Barium}-- A metallic element (atomic number 56, symbol Ba) whose sulfate salt (barium sulfate) is used as a radiocontrast agent due to its high atomic number, which effectively absorbs X-rays. Barium sulfate is insoluble and passes through the gastrointestinal tract without absorption, allowing radiographic visualization of the esophagus (barium swallow), stomach and duodenum (upper GI series), small bowel (barium follow-through), and colon (barium enema). Adverse effects include constipation and, rarely, barium impaction. Barium should not be used if perforation is suspected.

\medterm{Barium Enema}-- Fluoroscopic imaging of the large intestine after retrograde filling with barium contrast. Double-contrast (barium + air) provides better mucosal detail. Indications: colorectal cancer screening (historical, now CT colonography preferred), evaluation of colonic obstruction, diverticulosis, and postsurgical anatomy. Findings: filling defects (polyp, cancer), strictures, diverticula, and fistulas. Contraindications: suspected perforation (use water-soluble contrast). Largely replaced by colonoscopy.

\medterm{Barium Swallow}-- Fluoroscopic examination of the esophagus (pharynx to gastroesophageal junction) with barium contrast. Modified barium swallow: speech therapist evaluates oral and pharyngeal phases of swallowing. Indications: dysphagia, esophageal dysmotility (achalasia, diffuse esophageal spasm), hiatal hernia, GERD, stricture, and esophageal cancer. Findings: aspiration, cricopharyngeal bar, webs, rings (Schatzki), strictures, and diverticula (Zenker, epiphrenic).

\medterm{Basic Metabolic Panel}-- A panel of blood tests measuring serum glucose, sodium, potassium, chloride, bicarbonate, BUN, and creatinine. Also includes calcium in some institutions. Uses: assess hydration status, renal function, acid-base balance, and glucose control. Abnormalities suggest electrolyte disturbances, renal impairment, diabetes, or acid-base disorders. BMP is used for monitoring hospitalized patients and annual health screening.

\medterm{BCG Vaccine}-- Bacille Calmette-Guerin, a live attenuated vaccine derived from Mycobacterium bovis, used for TB prevention. Efficacy: 60-80\% against disseminated TB and TB meningitis in children; variable efficacy against pulmonary TB. Indications: part of routine childhood immunization in high-TB-burden countries. Not routinely used in low-risk countries (interferes with TST/IGRA screening). Also used as intravesical immunotherapy for non-muscle invasive bladder cancer.

\medterm{Biophysical Profile (BPP)}-- Antepartum fetal assessment combining ultrasound and NST. Components (2 points each): NST reactivity (reactive), fetal breathing movements (\(\ge\)1 episode of 30 seconds in 30 minutes), fetal movement (\(\ge\)3 discrete movements in 30 minutes), fetal tone (\(\ge\)1 episode of extension/flexion), and amniotic fluid volume (AFI \(\ge\)5 cm). Score: 8-10 (normal), 6 (equivocal, deliver if term, repeat if preterm), 4 (suspect asphyxia, deliver), 0-2 (probable asphyxia, deliver). Modified BPP: NST + AFI.

\medterm{Biopsy}-- Sampling of tissue for histopathologic examination. Types: needle biopsy (core needle: breast, liver, kidney, prostate; fine needle: thyroid, lymph node), excisional biopsy (entire lesion removal), incisional biopsy (partial sampling), endoscopic biopsy (GI, bronchoscopic), punch biopsy (skin), shave biopsy (skin), and stereotactic biopsy (breast guided by mammography). Essential for distinguishing benign from malignant, staging, and identifying inflammatory or infectious pathology.

\medterm{Blood Alcohol Level}-- Measurement of ethanol concentration in blood. Legal limit for driving: 0.08 g/dL (80 mg/dL) in most US states. Impaired: 0.04-0.08 g/dL. Severe intoxication: 0.15-0.30 g/dL. Potentially fatal: >0.30 g/dL. With chronic alcoholism: marked tolerance (may appear functional at 0.30+). Elimination rate: 15-25 mg/dL/hour. Uses: legal (DUI assessment), medical (alcohol intoxication, withdrawal management), and screening. Serum: whole blood ratio approximately 1.18:1.

\medterm{Blood Culture}-- Inoculation of blood into culture media to detect bacteremia or fungemia. Two sets (two bottles each: aerobic + anaerobic) from separate venipuncture sites, 20-30 mL total. Indications: fever of unknown origin, suspected sepsis, infective endocarditis, and immunocompromised with fever. Positive cultures with identification and susceptibility (MIC) guide targeted antibiotic therapy. Negative cultures do not rule out infection (especially with prior antibiotics, fastidious organisms, or intracellular pathogens).

\medterm{Blood Gas (Venous)}-- Sampling of venous blood (typically from a peripheral vein or central line) for pH, PvCO2, HCO3, and base excess. Less painful than ABG, avoids arterial puncture risks. Limitations: no PaO2 or SpO2 data; peripheral venous gases in shock may not reflect arterial pH. VBG pH is typically 0.03-0.05 lower than ABG; PvCO2 is 3-8 mmHg higher. Used for acid-base assessment and CO2 monitoring when oxygenation is not the primary concern.

\medterm{Bone Marrow Aspiration and Biopsy}-- Sampling of bone marrow, typically from the posterior iliac crest. Aspiration: liquid marrow for cytology (differential, morphology), flow cytometry, cytogenetics, and MRD. Biopsy: core of bone for cellularity, architecture, fibrosis, and infiltrative disease (lymphoma, metastatic cancer, granulomas). Indications: unexplained cytopenias, leukemia, lymphoma, multiple myeloma, myelodysplasia, myelofibrosis, and staging of certain malignancies. Contraindication: severe coagulopathy.

\medterm{Bone Scan (Nuclear Scintigraphy)}-- Nuclear medicine imaging using Tc-99m MDP (methylene diphosphonate) that accumulates in areas of increased bone turnover. Three phases: perfusion (immediate), blood pool (5 min), and delayed (2-3 hours). Indications: occult fracture, metastatic disease (breast, prostate, lung), osteomyelitis, complex regional pain syndrome (CRPS), avascular necrosis, and Paget disease. Findings: increased uptake in osteoblastic activity. Limitations: poor specificity (cannot distinguish tumor from infection or fracture). Used with SPECT for better localization.

\medterm{Brain Natriuretic Peptide (BNP)}-- A cardiac neurohormone released from the ventricles in response to volume overload and wall stress. BNP <100 pg/mL makes heart failure unlikely; >400 pg/mL suggests heart failure. NT-proBNP: <125 pg/mL (<75 years) or <450 pg/mL (>75 years) excluding HF. Also prognostic: rising levels indicate worsening HF and higher mortality. Elevated in HF (systolic and diastolic), pulmonary embolism, CKD, and sepsis. Used for diagnosis, severity assessment, and monitoring therapy.

\medterm{Breast MRI}-- Magnetic resonance imaging of the breast with IV gadolinium contrast. Indications: screening high-risk women (BRCA mutation, lifetime risk >20\%, prior chest radiation), evaluating extent of disease after breast cancer diagnosis, assessing response to neoadjuvant chemotherapy, implant evaluation, and problem-solving for equivocal mammogram/ultrasound. High sensitivity (>90\%) but lower specificity, leading to false positives. BI-RADS lexicon: kinetics (wash-in, washout) and morphology.

\medterm{Breast Ultrasound}-- Ultrasound imaging of the breast. Indications: evaluating palpable masses (distinguishing solid vs. cystic), characterizing mammographic abnormalities, guiding biopsy (core needle, fine needle, or vacuum-assisted), evaluating implant integrity, and imaging young/dense breast tissue. BI-RADS US classification: simple cyst (benign), complicated cyst, complex cystic and solid, and solid mass (assessed by shape, margins, orientation, echogenicity, and posterior features). Used with mammography for comprehensive breast evaluation.

\medterm{Bronchoscopy}-- Endoscopic examination of the tracheobronchial tree. Flexible bronchoscopy: most common, for diagnosis (lung mass, persistent cough, hemoptysis, suspected infection, ILD) and therapy (foreign body removal, bronchial washing, biopsy, BAL). Rigid bronchoscopy: for massive hemoptysis, large foreign bodies, and stent placement. Samples: BAL (bronchoalveolar lavage: cell count, culture, cytology), transbronchial biopsy (lung parenchyma), endobronchial biopsy (visible lesions), and EBUS-TBNA (mediastinal lymph node sampling).

\medterm{CA-125}-- Cancer antigen 125, a tumor marker most commonly used for ovarian cancer. Elevated in 80\% of epithelial ovarian cancers (less in mucinous). Non-specific: elevated in endometriosis, PID, pregnancy, fibroids, pancreatitis, cirrhosis, and other malignancies (pancreatic, breast, lung, endometrial). Normal: <35 U/mL. Uses: monitoring treatment response in known ovarian cancer, detecting recurrence, and differential diagnosis of pelvic mass (with Risk of Malignancy Index). Not recommended for screening (low specificity, high false positives).

\medterm{CA 19-9}-- A tumor marker for pancreatic cancer. Normal: <37 U/mL. Sensitivity: 70-80\% for pancreatic cancer (lower for early-stage). Specificity: limited (elevated in cholangitis, pancreatitis, cholecystitis, biliary obstruction, cirrhosis, and other GI malignancies). Uses: monitoring treatment response in known pancreatic cancer, detecting recurrence, and prognosis. Not recommended for screening (low sensitivity/specificity). Non-secretor phenotype (Lewis a-b-, 5-10\% of population) cannot produce CA 19-9, limiting its use.

\medterm{Capsule Endoscopy}-- A wireless, ingestible camera capsule (11 x 26 mm) that takes 50,000+ images as it traverses the GI tract. Used to visualize the small intestine (beyond reach of standard endoscopy). Indications: obscure GI bleeding, Crohn disease (small bowel involvement), suspected small bowel tumors, and surveillance in familial polyposis. Images transmitted to a data recorder worn on a belt, downloaded for review (3-4 hour study). Limitations: cannot biopsy or treat, capsule retention (1-2\% risk, higher in Crohn strictures), and poor visualization of the stomach and colon. PillCam (Given Imaging) and EndoCapsule (Olympus) are common systems.

\medterm{Cardiac Catheterization}-- Invasive procedure to evaluate cardiac anatomy, hemodynamics, and coronary arteries. Right heart catheterization: Swan-Ganz catheter measures RA, RV, PA pressures, PCWP, cardiac output (thermodilution), and SVR/PVR. Left heart catheterization: coronary angiography and ventriculography. Indications: CAD evaluation (stable angina with high-risk features, positive stress test, acute coronary syndromes), valvular disease (assess severity), cardiomyopathy, and pre-transplant evaluation. Complications: hematoma, bleeding, CVA, MI, arrhythmia.

\medterm{Cardiac Stress Test}-- Evaluation of cardiac function and ischemia under increased cardiac workload. Exercise stress test (treadmill, Bruce protocol): ECG monitoring for ST-segment changes, symptoms, BP response, and arrhythmias. Sensitivity 65-75\% for CAD. Pharmacologic stress test: dobutamine (inotrope/chronotrope), regadenoson/adenosine/dipyridamole (vasodilator) — used for patients unable to exercise. Imaging: stress echo (wall motion), nuclear perfusion (SPECT, PET), and stress MRI.

\medterm{Cardiopulmonary Exercise Testing (CPET)}-- Comprehensive exercise test with breath-by-breath analysis of respiratory gases (VO2, VCO2, VE). Indications: unexplained dyspnea, pre-operative risk assessment (lung resection, cardiac surgery), evaluation of exercise tolerance in heart failure (peak VO2 predicts prognosis and guides transplant listing), and differentiating cardiac vs. pulmonary limitation. Key parameters: peak VO2 (normal >80\% predicted), anaerobic threshold (AT), VEVCO2 slope (<30 normal, >34 poor prognosis in HF), and oxygen pulse (VO2/HR). Maximum exercise test with 12-lead ECG and BP monitoring.

\medterm{Cardioversion (Synchronized)}-- Delivery of a low-energy electrical shock synchronized to the QRS complex to terminate arrhythmias. Synchronization: shock delivered on the R wave (not T wave, which induces VF). Energy: atrial fibrillation (100-200 J biphasic), atrial flutter (50-100 J), SVT (50-100 J), stable VT (100 J). Success: 70-90\% for AF, >95\% for atrial flutter. Pre-treatment: anticoagulation for >3 weeks or TEE-guided to exclude LA thrombus (for AF >48 hours). Complications: thromboembolism, skin burns, arrhythmias, and sedation risks.

\medterm{Caries Risk Assessment}-- Systematic evaluation of an individual's likelihood of developing dental caries. Low risk: no active caries, good oral hygiene, regular fluoride exposure. Moderate risk: one or more active lesions in past 3 years, fair oral hygiene. High risk: multiple active lesions, poor oral hygiene, frequent sugar intake, low fluoride, xerostomia, smoking, or special healthcare needs. Based on clinical exam (caries, restorations), radiographs (incipient lesions), plaque assessment, and dietary/medical history. Guides recall interval (3-12 months) and preventive interventions.

\medterm{Cariogram}-- A computer-based program that illustrates the interaction of caries-related factors and provides an overall caries risk assessment. Factors: diet (content, frequency), bacteria (mutans streptococci, lactobacilli), fluoride exposure, saliva (flow rate, buffering capacity), plaque amount, and clinical/social history. Output: pie chart showing the relative weight of different etiological factors and the probability of avoiding new caries. Used for patient education and preventive planning.

\medterm{Carotid Duplex Ultrasound}-- Ultrasound evaluation of the carotid arteries for atherosclerotic plaque and stenosis. Measures peak systolic velocity (PSV) and end-diastolic velocity (EDV) in the internal carotid artery (ICA). Criteria: normal (<50\% stenosis: PSV <125 cm/s), 50-69\% (PSV 125-230, EDV 40-100), 70-99\% (PSV >230, EDV >100), and occlusion (no detectable flow). Indications: TIA/stroke evaluation, asymptomatic bruit, pre-CABG screening, and surveillance of known stenosis. Plaque morphology: echolucent (high risk) vs. calcified.

\medterm{Carpal Tunnel Syndrome Tests}-- Provocative tests for diagnosing carpal tunnel syndrome. Tinel sign: percuss over median nerve at the carpal tunnel $\rightarrow$ tingling in median distribution. Phalen sign: hold wrists in full flexion for 60 seconds $\rightarrow$ reproduces symptoms. Durkan test (carpal compression): direct compression over carpal tunnel for 30 seconds $\rightarrow$ paresthesias. Thenar atrophy suggests advanced disease. Confirmatory: nerve conduction studies show prolonged distal latencies and decreased conduction velocity across the carpal tunnel.

\medterm{Catecholamines (Urine/Plasma)}-- Measurement of catecholamines (epinephrine, norepinephrine, dopamine) and their metabolites (metanephrines, VMA) for diagnosis of pheochromocytoma and paraganglioma. Plasma free metanephrines: sensitivity 96-100\%. 24-hour urinary fractionated metanephrines: sensitivity 90-95\%. Elevated levels require tumor localization (CT/MRI, MIBG scan). False positives: stress, exercise, smoking, tricyclic antidepressants, MAOIs, and levodopa.

\medterm{CD4 Count}-- Absolute count of CD4+ T-helper lymphocytes, essential for monitoring HIV disease progression and treatment response. Normal: 500-1500 cells/uL. HIV staging: >500 (asymptomatic, start ART regardless), 200-500 (consider prophylaxis for some OIs), <200 (AIDS-defining: start PCP prophylaxis, TB screening, high risk of OIs). <100: risk of CMV, MAC, toxoplasmosis, cryptococcosis. Used with HIV viral load (HIV RNA) to guide ART initiation and monitor response (goal: undetectable viral load, CD4 recovery).

\medterm{CEA (Carcinoembryonic Antigen)}-- A tumor marker primarily for colorectal cancer. Normal: <3 ng/mL (non-smokers), <5 ng/mL (smokers). Uses: monitoring treatment response in metastatic colorectal cancer, surveillance for recurrence after curative resection (q3-6 months for 2 years, then q6-12 months for 5 years), and prognosis (preoperative level). Non-specific elevations: smoking, pancreatitis, IBD, cirrhosis, hepatitis, and other malignancies (lung, gastric, pancreatic, breast). Not recommended for screening.

\medterm{Central Line Placement}-- Insertion of a catheter into a central vein (internal jugular, subclavian, or femoral), with the tip in the superior (or inferior) vena cava. Ultrasound guidance is standard. Indications: administration of vasopressors, TPN, chemotherapy, long-term antibiotics, central venous pressure monitoring, hemodialysis, and plasmapheresis. Complications: arterial puncture, pneumothorax (subclavian > IJ), hemothorax, air embolism, catheter-related bloodstream infection, thrombosis, and malposition. Confirm position with chest X-ray (tip at SVC-RA junction, 2-3 cm above RA). Suture and sterile dressing.

\medterm{Cerebrospinal Fluid Analysis}-- Laboratory examination of CSF obtained via lumbar puncture. Indications: meningitis, encephalitis, subarachnoid hemorrhage, demyelinating disease (MS: oligoclonal bands, elevated IgG index), and carcinomatous meningitis. Components: opening pressure (normal 10-20 cmH2O), appearance (clear vs. cloudy vs. bloody), cell count/differential (normal <5 WBC, all monocytes; neutrophils in bacterial), glucose (CSF:serum ratio >0.6), protein (normal 15-45 mg/dL), Gram stain, culture, and PCR.

\medterm{Cervical Screening (Pap Smear)}-- Screening for cervical cancer by collecting cells from the transformation zone of the cervix. Liquid-based cytology (ThinPrep) or conventional Pap smear. HPV co-testing (high-risk HPV 16, 18, others) improves sensitivity. Screening guidelines: start at age 21; ages 21-29: Pap q3 years; ages 30-65: Pap + HPV q5 years or Pap q3 years; over 65: discontinue if adequate prior negative screening. Results: ASCUS, LSIL, HSIL, AGC, and carcinoma. Management per ASCCP risk-based guidelines.

\medterm{Chest Tube Thoracostomy}-- Insertion of a tube into the pleural space to drain air, fluid, or blood. Indications: pneumothorax (large, tension, or on mechanical ventilation), hemothorax, pleural effusion (empyema, malignant, transudative), and postoperative (cardiac, thoracic surgery). Tube sizes: small (8-14F for air), medium (16-24F for effusion), large (28-40F for hemothorax/empyema). Placement: triangle of safety (4th-5th intercostal space, midaxillary line). Complications: bleeding, infection, lung laceration, and re-expansion pulmonary edema.

\medterm{Chorionic Villus Sampling}-- Prenatal diagnostic test at 10-13 weeks gestation, earlier than amniocentesis. Transabdominal or transcervical biopsy of chorionic villi (placental tissue). Indications: advanced maternal age, abnormal first-trimester screen, carrier screening, or family history of genetic disorders. Analyzed for karyotype, FISH, and microarray. Advantages: earlier diagnosis than amniocentesis (allows earlier termination if desired). Risks: miscarriage (0.5-1\%, slightly higher than amniocentesis) and limb defects if performed <10 weeks. Does not detect neural tube defects (requires AFP at 16-18 weeks).

\medterm{Coagulation Profile}-- Panel of tests assessing the coagulation cascade. PT (prothrombin time, 11-14 sec): measures extrinsic and common pathways; elevated in vitamin K deficiency, warfarin, liver disease, DIC. INR standardizes PT for warfarin monitoring. PTT (partial thromboplastin time, 25-35 sec): measures intrinsic and common pathways; elevated in hemophilia A/B, heparin, lupus anticoagulant, factor deficiencies. Fibrinogen (200-400 mg/dL): acute phase reactant; low in DIC, liver disease. Mixing studies: corrects in factor deficiency, does not correct in inhibitor.

\medterm{Colonoscopy}-- Endoscopic examination of the entire colon from cecum to rectum. Indications: colorectal cancer screening (starting age 45 for average risk), surveillance (polyps: q3-5y; cancer: q1y; IBD: q1-2y for dysplasia), diagnostic (hematochezia, chronic diarrhea, iron deficiency anemia, change in bowel habits), and therapeutic (polypectomy, hemostasis, stricture dilation, stent placement). Requires full bowel preparation (polyethylene glycol, split-dose). Complications: perforation (0.03-0.3\%), bleeding (0.1-1\%), and sedation-related.

\medterm{Colorectal Cancer Screening}-- See Colonoscopy, FIT, FIT-DNA, CT Colonography.

\medterm{Colorectal Cancer Screening Tests}-- Multiple screening modalities. Stool-based: fecal immunochemical test (FIT, annual), high-sensitivity guaiac FOBT (annual), and multitarget stool DNA (FIT-DNA/Cologuard, q3y). Direct visualization: colonoscopy (q10y preferred), CT colonography (q5y), flexible sigmoidoscopy (q5-10y), and capsule colonoscopy (q5y). Guidelines: average risk start at age 45; stop at age 75 if adequate prior screening and negative; age 76-85: individual decision based on life expectancy and comorbidities. FIT is the most cost-effective option.

\medterm{Colposcopy}-- Magnified visualization of the cervix, vagina, and vulva after application of acetic acid. Indication: abnormal Pap smear (ASCUS with high-risk HPV, LSIL, HSIL, AGC). Acetic acid causes acetowhite lesions in dysplastic epithelium. Lugol iodine (Schiller test): normal glycogenated epithelium stains brown; abnormal epithelium does not. Targeted biopsy (punch biopsies and endocervical curettage) of abnormal areas to confirm CIN (cervical intraepithelial neoplasia). Grading: CIN 1 (mild), CIN 2 (moderate), CIN 3 (severe/carcinoma in situ).

\medterm{Complete Blood Count (CBC)}-- A panel of blood tests evaluating cellular components: WBC (leukocytes: normal 4.5-11.0 K/uL), RBC (erythrocytes: male 4.5-5.9, female 4.1-5.1 M/uL), hemoglobin (male 13.5-17.5, female 12.0-16.0 g/dL), hematocrit (male 41-53\%, female 36-46\%), platelets (150-450 K/uL), RBC indices (MCV, MCH, MCHC), and differential (neutrophils, lymphocytes, monocytes, eosinophils, basophils). Uses: anemia, infection, bleeding, inflammation, and hematologic malignancy screening.

\medterm{Comprehensive Metabolic Panel}-- Extended version of BMP including albumin, total protein, alkaline phosphatase, ALT, AST, and bilirubin (total and direct). Evaluates liver function, kidney function, electrolytes, glucose, calcium, and protein status. Uses: preoperative assessment, medication monitoring, and liver/kidney disease evaluation. Flags acute and chronic conditions: transaminitis (hepatitis, NAFLD), elevated ALP (obstruction, bone disease), low albumin (cirrhosis, malnutrition).

\medterm{Coronary Calcium Score}-- Non-contrast CT measurement of coronary artery calcium (Agatston score). Score 0: no plaque, very low event risk (near 0\% over 5 years). 1-100: mild plaque, moderate risk. 101-400: moderate plaque, increased risk. >400: extensive plaque, high risk. Used in ASCVD risk assessment for intermediate-risk patients when risk-based decision is uncertain. Independent predictor of cardiovascular events beyond traditional risk factors. Zero calcium score is an excellent negative predictor.

\medterm{COVID-19 Test (SARS-CoV-2)}-- Detection of SARS-CoV-2. NAAT (RT-PCR): gold standard, nasopharyngeal/nasal swab, highly sensitive/specific, results in 1-48 hours. Rapid antigen test: lateral flow, results in 15 minutes, sensitive in high viral load (early symptomatic), less sensitive in asymptomatic. Serology: anti-S and anti-N antibodies, indicates prior infection or vaccination, not used for acute diagnosis. Used for diagnosis, infection control, and public health surveillance. Antiviral therapy (nirmatrelvir/ritonavir, remdesivir) for high-risk patients.

\medterm{C-Reactive Protein (CRP)}-- An acute-phase reactant produced by the liver in response to inflammation (IL-6). CRP rises within 4-6 hours of inflammatory stimulus and peaks at 36-50 hours. Uses: monitoring inflammatory disease activity (RA, vasculitis), detecting infection postoperatively, and cardiovascular risk stratification (hs-CRP: high-sensitivity). CRP >10 mg/L suggests significant inflammation or infection. More sensitive than ESR for acute changes.

\medterm{Creatine Kinase (CK)}-- An enzyme found in muscle, brain, and heart. CK-MB: cardiac isoenzyme (used historically for MI diagnosis, largely replaced by troponin). CK-MM: skeletal muscle (elevated in rhabdomyolysis, trauma, statins, myopathy). Total CK: elevated in muscle injury, myocardial infarction, seizures, hypothyroidism, and statin myopathy. Strikingly elevated in rhabdomyolysis (>5000-100,000 U/L). Used to monitor rhabdomyolysis (risk of AKI) and inflammatory myopathies.

\medterm{Cryoglobulins}-- Immunoglobulins that precipitate at cold temperatures (<37°C) and redissolve on warming. Types: Type I (monoclonal IgM, in MGUS, Waldenstrom, myeloma), Type II (mixed: monoclonal + polyclonal, in HCV infection), Type III (mixed polyclonal, in autoimmune disease). Presents with palpable purpura, arthralgias, neuropathy, and glomerulonephritis (membranoproliferative). Diagnosis: serum collected and transported at 37°C, cold-precipitated, measured as cryocrit. Treatment: address underlying cause (HCV: DAAs), immunosuppression.

\medterm{CT Angiography (CTA)}-- CT imaging with IV contrast timed to visualize arteries. Coronary CTA: assesses coronary artery stenosis and plaque, calcium scoring (Agatston). CTA chest: pulmonary embolism (PE protocol), aortic dissection. CTA head/neck: stroke evaluation, aneurysm, dissection, and vasculitis. CTA abdomen: mesenteric ischemia, renal artery stenosis, and AAA. 64-slice or better scanners required. Advantages: non-invasive, three-dimensional reconstruction, and rapid acquisition. Disadvantage: radiation and contrast exposure.

\medterm{CT Colonography (Virtual Colonoscopy)}-- CT imaging of the colon after bowel preparation and air insufflation. Indications: colorectal cancer screening (alternative to colonoscopy), incomplete colonoscopy due to obstructing mass, and patients unable to undergo colonoscopy. Sensitivity >90\% for polyps >10mm. Requires same bowel prep as colonoscopy. Limitations: cannot biopsy or remove polyps; positive findings require optical colonoscopy. No sedation required.

\medterm{CT Scan (Computed Tomography)}-- Cross-sectional imaging using X-rays and computer reconstruction. Non-contrast: evaluate hemorrhage, stones, and bone. IV contrast (iodinated): evaluate inflammation, infection, tumor, and vascular structures. Oral contrast: opacify bowel. Window settings: bone (fractures, calcifications), soft tissue (organs, masses), lung (pulmonary parenchyma), and brain (gray-white differentiation). Radiation dose varies (2-8 mSv for head, 5-15 mSv for body). Hounsfield units quantify tissue density.

\medterm{Culture and Sensitivity}-- Microbiological culture of a clinical specimen (blood, urine, sputum, wound swab, CSF) to identify pathogens and determine antimicrobial susceptibility. Culture media: blood agar (most bacteria), MacConkey (Gram-negative), chocolate agar (fastidious: N. gonorrhoeae, H. influenzae), and selective media. Automated systems (VITEK, MALDI-TOF) speed identification. MIC (minimum inhibitory concentration) reported as susceptible, intermediate, or resistant per CLSI breakpoints. Guides targeted antimicrobial therapy.

\medterm{Cystoscopy}-- Endoscopic examination of the urethra, bladder, and ureteral orifices. Flexible cystoscopy: office-based, local anesthesia. Rigid cystoscopy: OR, allows larger instruments for interventions. Indications: hematuria (gross or microscopic), recurrent UTIs, bladder tumor surveillance, lower urinary tract symptoms (obstruction), ureteral stent placement/removal, and foreign body removal. Findings: tumors, stones, strictures, trabeculation, diverticula, and inflammation. Complications: infection, hematuria, and perforation.

\medterm{Cytogenetics (Karyotype)}-- Visualization of chromosomes under light microscopy from dividing cells (bone marrow, blood, amniotic fluid, tissue). Staining: G-banding (Giemsa, light and dark bands). Resolution: 400-550 bands per haploid genome. Detects: aneuploidy (trisomy 21, 18, 13, Turner, Klinefelter), structural abnormalities (translocations, deletions, inversions, duplications), and marker chromosomes. Essential in: prenatal diagnosis, hematologic malignancy (t(9;22) BCR-ABL in CML, t(15;17) PML-RARA in APL, t(8;14) MYC in Burkitt), and constitutional genetic disorders.

\medterm{D-Dimer}-- A fibrin degradation product, elevated in thrombosis (PE, DVT, DIC). High sensitivity but low specificity. D-dimer <500 ng/mL (ELISA) has high negative predictive value for VTE; used to rule out PE/DVT in patients with low or moderate pretest probability (Wells criteria). Age-adjusted D-dimer (>50 years: age x 0.1 mg/L) improves specificity in older patients. Also elevated in infection, inflammation, surgery, trauma, pregnancy, and malignancy.

\medterm{Deep Vein Thrombosis Ultrasound}-- Venous duplex ultrasound for DVT diagnosis. Complete compression ultrasound (CUS): venous segments are compressed; non-compressible = DVT (96\% sensitivity for proximal DVT). Whole-leg ultrasound (proximal + distal): for symptomatic patients. Findings: intraluminal thrombus, absent flow, lack of compressibility, and abnormal Doppler signals. Distal DVT (calf) is less clinically significant; management varies.

\medterm{Dental Radiographs}-- Imaging for dental diagnosis. Intraoral: periapical (tooth root and surrounding bone), bitewing (interproximal caries detection, crestal bone level), and occlusal (larger area, localization). Panoramic (extraoral): wide view of jaws, teeth, sinuses, and TMJ; used for impacted teeth, fractures, pathology (cysts, tumors), and implant planning. CBCT (cone-beam CT): 3D imaging for implant planning and complex surgical cases. Frequency based on caries risk.

\medterm{Direct Coombs Test (DAT)}-- Direct antiglobulin test: detects antibodies or complement (C3) bound to the surface of RBCs in vivo. Positive in: autoimmune hemolytic anemia (warm: IgG, cold: C3), drug-induced hemolytic anemia, hemolytic transfusion reaction, and hemolytic disease of the newborn (HDN). Grading: 1+ to 4+. Used with indirect Coombs test (antibody screen for transfusion).

\medterm{Doppler Ultrasound}-- Ultrasound technique using the Doppler effect to evaluate blood flow. Color Doppler: direction and velocity of flow (red = toward transducer, blue = away). Spectral Doppler: waveform analysis (arterial: pulsatile; venous: continuous). Pulsed-wave: measures flow at a specific depth. Continuous-wave: measures high velocities (valvular stenosis). Uses: assessment of DVT, carotid stenosis, PAD, renal artery stenosis, valvular heart disease, and pregnancy (fetal umbilical artery, middle cerebral artery).

\medterm{DXA Scan (Dual-Energy X-ray Absorptiometry)}-- The gold standard for measuring bone mineral density (BMD) and diagnosing osteoporosis. Measures at the lumbar spine, total hip, and femoral neck. T-score: normal (\(\ge\)-1.0), osteopenia (-1.0 to -2.5), osteoporosis (\(\le\)-2.5). Z-score (compared to age-matched) used in premenopausal women and men <50. Indications: women \(\ge\)65, men \(\ge\)70, younger adults with risk factors (prior fragility fracture, steroid use, hypogonadism, malabsorption). FRAX calculates 10-year fracture risk. Used for monitoring treatment response (repeat q2y).

\medterm{Echocardiography}-- Ultrasound imaging of the heart. Transthoracic (TTE): standard, assesses LV/RV function, wall motion, valves (stenosis, regurgitation), pericardium (effusion), and chamber sizes. Transesophageal (TEE): higher resolution of valves (mitral, aortic), left atrial appendage, and for endocarditis evaluation. Stress echo: images before and after stress (exercise or dobutamine) for inducible wall motion abnormalities (ischemia). Doppler: measures blood flow velocities across valves (pressure gradients) and estimates pulmonary artery pressure.

\medterm{Electrocardiogram (ECG / EKG)}-- A recording of the heart's electrical activity from 12 leads (I, II, III, aVR, aVL, aVF, V1-V6). P wave: atrial depolarization. PR interval (120-200 ms): conduction through AV node. QRS (80-110 ms): ventricular depolarization. ST segment: early repolarization. T wave: ventricular repolarization. QT interval: corrected for heart rate (QTc, normal <450 ms in men, <460 ms in women). Uses: chest pain (ischemia/infarction), arrhythmia diagnosis, syncope, palpitations, and pre-surgical evaluation.

\medterm{Electroencephalogram (EEG)}-- Recording of the brain's electrical activity from scalp electrodes. Background rhythms: alpha (8-13 Hz, relaxed, eyes closed, occipital), beta (13-30 Hz, alert/active), theta (4-8 Hz, drowsy/light sleep, pathological in awake adults), delta (<4 Hz, deep sleep, pathological in awake). Uses: seizure diagnosis and classification, epilepsy monitoring, altered mental status, brain death determination, and sleep disorders. Activation: hyperventilation, photic stimulation, and sleep deprivation.

\medterm{Electrophysiology Study (EPS)}-- Invasive cardiac catheterization with multielectrode catheters placed in the heart to map the cardiac conduction system. Measures: sinus node function (SNRT, SACT), AV node function (AH interval, HV interval), and inducibility of arrhythmias (programmed electrical stimulation). Indications: evaluation of syncope, sustained palpitations, wide-complex tachycardia, supraventricular tachycardia (AVNRT, AVRT, atrial flutter), ventricular tachycardia, and risk stratification in Brugada syndrome or hypertrophic cardiomyopathy. Often followed by catheter ablation.

\medterm{Endoscopic Retrograde Cholangiopancreatography}-- Endoscopic cannulation of the ampulla of Vater for retrograde injection of contrast into the biliary and pancreatic ducts. Combined diagnostic and therapeutic: ERCP with sphincterotomy, stone extraction, stent placement, and brush cytology. Indications: choledocholithiasis (common bile duct stones), malignant biliary obstruction (palliative stenting), biliary stricture, and sphincter of Oddi dysfunction. Risks: pancreatitis (5-10\%), cholangitis, perforation, and bleeding post-sphincterotomy.

\medterm{Endoscopic Ultrasound (EUS)}-- Combined endoscopy and ultrasound for high-resolution imaging of the GI wall and adjacent structures. Linear echoendoscope: allows FNA of mediastinal, pancreatic, and peripancreatic lesions. Radial echoendoscope: 360-degree view for staging GI cancers (TNM staging). Indications: pancreatic cancer staging, mediastinal lymph node sampling (EBUS-TBNA or EUS-FNA), subepithelial GI lesions, choledocholithiasis, and celiac plexus neurolysis for pain. Complications: perforation (0.03\%), pancreatitis (from pancreatic FNA), and infection.

\medterm{Erythrocyte Sedimentation Rate (ESR)}-- A non-specific inflammatory marker measuring the rate at which erythrocytes settle in 1 hour. Normals: males age/2, females (age+10)/2. Markedly elevated >100 mm/hr: suggests infection (endocarditis), autoimmune (giant cell arteritis, polymyalgia rheumatica, SLE, RA), malignancy (multiple myeloma), and chronic kidney disease. Faster in anemia (rouleaux formation). Slower in polycythemia and spherocytosis. Uses: monitoring disease activity in inflammatory conditions. Less sensitive than CRP for acute changes.

\medterm{Esophageal Manometry}-- Measurement of esophageal pressures at rest and during swallows. Components: upper esophageal sphincter (UES), esophageal body (amplitude, peristalsis), and lower esophageal sphincter (LES, resting pressure and relaxation). High-resolution manometry (HRM) with Chicago Classification: diagnoses achalasia (incomplete LES relaxation, absent peristalsis), esophagogastric junction outflow obstruction, distal esophageal spasm, hypercontractile (jackhammer) esophagus, and ineffective esophageal motility. Indications: dysphagia, non-cardiac chest pain, and pre-operative evaluation for anti-reflux surgery.

\medterm{Esophageal pH Monitoring}-- 24-hour ambulatory monitoring of esophageal acid exposure. Catheter-based (transnasal pH probe placed 5 cm above LES) or wireless (Bravo capsule endoscopically attached to esophageal mucosa). Measures: \% time pH <4 (total, upright, supine), DeMeester score (>14.7 abnormal), and symptom-reflux correlation. Indications: refractory GERD, atypical GERD symptoms (chest pain, cough, asthma), pre-operative evaluation for anti-reflux surgery, and monitoring of therapy. Off acid suppression for diagnostic testing.

\medterm{Esophagogastroduodenoscopy (EGD)}-- Upper endoscopy: flexible endoscopic examination of the esophagus, stomach, and duodenum. Indications: dysphagia, GERD (surveillance of Barrett esophagus), dyspepsia (alarm symptoms: age >60, weight loss, anemia, dysphagia, bleeding), upper GI bleeding (variceal and non-variceal: ulcers, Mallory-Weiss tears), and surveillance for polyps (FAP, gastric intestinal metaplasia). Therapeutic: hemostasis (cautery, clips, sclerotherapy), polypectomy, dilation (strictures), and PEG tube placement.

\medterm{Evoked Potentials (VEP, BAER, SSEP)}-- Electrophysiologic tests measuring the CNS electrical response to specific sensory stimuli. Visual evoked potentials (VEP): checkboard pattern reversal, measures P100 latency; delayed in optic neuritis/MS. Brainstem auditory evoked responses (BAER): click stimuli, waves I-V; delayed in acoustic neuroma, MS, brainstem stroke. Somatosensory evoked potentials (SSEP): electrical stimulation of peripheral nerves (median, tibial); used intraoperatively for spinal cord monitoring (scoliosis surgery) and prognostic in spinal cord injury and hypoxic-ischemic encephalopathy.

\medterm{Exercise Stress Test}-- Treadmill exercise with ECG monitoring using the Bruce protocol (3-minute stages increasing speed and grade). Indications: diagnosis of CAD (chest pain evaluation), functional capacity assessment, and prognosis after MI. Contraindications: acute MI (within 2 days), unstable angina, decompensated heart failure, severe aortic stenosis, and uncontrolled arrhythmias. Positive test: \(\ge\)1 mm horizontal or downsloping ST depression at 80 ms after the J point. Duke Treadmill Score (DTS) provides prognostic information (duration, ST deviation, angina).

\medterm{Extraocular Muscle Function Testing}-- Assessment of the six extraocular muscles innervated by CN III (medial rectus: adduction; superior rectus: elevation; inferior rectus: depression; inferior oblique: extorsion), CN IV (superior oblique: intorsion and depression), and CN VI (lateral rectus: abduction). Tested by asking the patient to follow a target in the six cardinal directions of gaze. Abnormalities: strabismus (misalignment), diplopia (double vision), and palsy. Cover-uncover test and alternate cover test identify phoria vs. tropia.

\medterm{Fecal Calprotectin}-- A protein released by neutrophils in the stool, a marker of intestinal inflammation. Normal: <50 mcg/g. Elevated in IBD (>150-200 mcg/g for active Crohn, ulcerative colitis). Distinguishes IBD from IBS (IBS: normal). Monitoring: correlates with endoscopic disease activity; reduction indicates treatment response. Also elevated in: NSAID enteropathy, infectious colitis, and colorectal cancer. False negatives: skip lesions in Crohn. Used to avoid unnecessary colonoscopies and guide treatment escalation.

\medterm{Fecal Immunochemical Test (FIT)}-- A stool-based screening test for colorectal cancer that detects human hemoglobin in a single stool sample using antibodies. Annual screening. Sensitivity: 74-80\% for cancer, 25-30\% for advanced adenomas. Specificity: 94-96\%. No dietary restrictions required. Positive test requires follow-up colonoscopy. More specific than guaiac FOBT. Cost-effective and non-invasive. Compared to colonoscopy: lower sensitivity but more accessible and less invasive.

\medterm{FeNO (Fractional Exhaled Nitric Oxide)}-- A non-invasive biomarker of type 2 airway inflammation (eosinophilic). Normal: <25 ppb (adults), <35 ppb (children). Elevated (>50 ppb): eosinophilic asthma, good response to inhaled corticosteroids. Low (<25 ppb): non-eosinophilic asthma, neutrophilic inflammation. Used to: diagnose asthma (when spirometry is equivocal), predict steroid responsiveness, monitor adherence to ICS, and guide treatment (adjust ICS dose to normalize FeNO). Affected by smoking (falsely low), atopy, and certain foods.

\medterm{Ferritin}-- Intracellular iron storage protein reflecting total body iron stores. Low ferritin: iron deficiency anemia (<30 ng/mL, absolute; <100 can be consistent with iron deficiency in inflammation). High ferritin: iron overload (hemochromatosis, repeated transfusions), inflammation (acute phase reactant), liver disease, malignancy, and hemophagocytic lymphohistiocytosis (HLH). Used with serum iron, TIBC, and transferrin saturation for complete iron studies. Transferrin saturation >45\% suggests iron overload.

\medterm{Fetal Ultrasound}-- Ultrasound imaging during pregnancy. First trimester (6-12 weeks): confirm pregnancy, gestational age (crown-rump length), number of fetuses, viability (heart activity), nuchal translucency (screening for aneuploidy at 11-13 weeks). Second trimester (18-22 weeks): anatomy survey (fetal head, chest, heart, abdomen, spine, extremities, placenta, amniotic fluid), gender determination. Third trimester: growth assessment (biometric parameters: BPD, HC, AC, FL), fetal weight estimation, placental location, amniotic fluid index, and biophysical profile. Doppler: umbilical artery, middle cerebral artery.

\medterm{FISH (Fluorescence In Situ Hybridization)}-- A molecular cytogenetic technique using fluorescent DNA probes that hybridize to specific chromosomal regions. Faster than karyotype (24-48 hours, vs. 7-14 days). Indications: rapid prenatal aneuploidy testing (trisomy 13, 18, 21, X, Y), detection of microdeletion syndromes (22q11.2 DiGeorge, TP53 in CLL), gene rearrangements in cancer (BCR-ABL, PML-RARA, MLL rearrangements, MYC, IGH), and HER2 amplification in breast cancer. Can be performed on interphase nuclei (does not require dividing cells). Higher resolution than karyotype.

\medterm{Flexible Sigmoidoscopy}-- Endoscopic examination of the rectum and sigmoid colon (up to 60 cm). Indications: colorectal cancer screening (q5-10 years, combined with FIT), evaluation of rectal bleeding, chronic diarrhea, and surveillance of known distal colonic disease (ulcerative colitis). Advantages over colonoscopy: no sedation required in most cases, less bowel prep (enema only), lower cost and risk. Limitations: cannot visualize the proximal colon; positive findings require full colonoscopy. Largely replaced by colonoscopy for screening in the US.

\medterm{Flow Cytometry}-- Technique that measures physical and chemical characteristics of cells or particles in a fluid stream, using laser-based detection. Commonly used in hematology: immunophenotyping of leukemia and lymphoma (CD markers: CD3, CD19, CD20, CD34, CD33, CD13), detection of minimal residual disease, and evaluation of immunodeficiency disorders. Analysis of DNA content (ploidy) in neoplasms. Also used in stem cell enumeration (CD34+) for transplant.

\medterm{Fluoroscopy}-- Real-time X-ray imaging used to visualize dynamic processes. Common applications: upper GI series (barium swallow, barium meal), lower GI series (barium enema), voiding cystourethrogram (VCUG for VUR), hysterosalpingography (tubal patency), angiography, and interventional procedures (ERCP, PICC line placement, joint injection, pacemaker implantation, and catheterization). Utilizes continuous or pulsed X-ray beam. Radiation exposure is a concern for both patients and operators.

\medterm{Genetic Testing (Sequencing)}-- Analysis of DNA to identify pathogenic variants associated with genetic disorders. Sanger sequencing: gold standard for single-gene testing (CFTR for cystic fibrosis, BRCA1/2 for hereditary breast/ovarian cancer). Next-generation sequencing (NGS): massively parallel sequencing enabling panel testing (multiple genes), exome sequencing (all exons, 1-2\% of genome), and genome sequencing (whole genome). Detects SNPs, indels, and copy number variants. Indications: confirm a suspected genetic diagnosis, carrier screening, prenatal diagnosis, pharmacogenomics, and cancer genomics. Variant interpretation: pathogenic, likely pathogenic, VUS (variant of uncertain significance), likely benign, benign.

\medterm{Glucose Tolerance Test}-- Testing for diabetes mellitus and gestational diabetes. 75g OGTT (oral glucose tolerance test): measure plasma glucose at 0, 60, and 120 minutes. Diagnostic: fasting \(\ge\)126 mg/dL, 2-hour \(\ge\)200 mg/dL. Gestational diabetes screening: 50g glucose challenge test (GCT) at 24-28 weeks; if >140 mg/dL at 1 hour, proceed to 100g 3-hour OGTT. Carpenter-Coustan criteria: \(\ge\)2 values elevated (fasting >95, 1h >180, 2h >155, 3h >140). Glucose tolerance categories: normal, impaired fasting glucose (IFG 100-125), impaired glucose tolerance (IGT 140-199), and diabetes.

\medterm{Gram Stain}-- A differential staining technique that classifies bacteria as Gram-positive (purple, thick peptidoglycan) or Gram-negative (pink/red, thin peptidoglycan + outer membrane). Also reports morphology (cocci, rods, spirals) and arrangement (clusters, chains, pairs). Used to guide empiric antibiotic therapy while awaiting culture. Quality: epithelial cells suggest contamination; polymorphonuclear leukocytes suggest true infection. Foundational in microbiology.

\medterm{Haptoglobin}-- An alpha-2 globulin that binds free hemoglobin, preventing oxidative damage and renal excretion. Low/undetectable: acute hemolysis (intravascular and extravascular), liver disease (decreased production), and hereditary deficiency. Elevated: acute phase reaction (inflammation, infection, malignancy). Used with LDH and bilirubin to confirm hemolytic anemia. Normal: 30-200 mg/dL.

\medterm{hCG (Human Chorionic Gonadotropin)}-- A glycoprotein hormone produced by the placenta, detected in urine and serum for pregnancy diagnosis. Qualitative: positive \(\ge\)25 mIU/mL (urine). Quantitative: serum hCG levels — <5 mIU/mL (negative), 5-25 (indeterminate), >25 (positive). In early pregnancy: hCG doubles every 48-72 hours until 8-10 weeks. Abnormal levels: ectopic pregnancy (slow rise), miscarriage (declining), or molar pregnancy (very high). Also a tumor marker for gestational trophoblastic disease and germ cell tumors.

\medterm{Hemoglobin A1c}-- Glycated hemoglobin reflecting average blood glucose over the preceding 2-3 months. Normal: <5.7\%. Prediabetes: 5.7-6.4\%. Diabetes: \(\ge\)6.5\% (diagnostic, confirmed with repeat). Targets: generally <7\% for most adults with diabetes (individualized: <6.5\% for some, <8\% for frail/elderly). Used for diabetes diagnosis and monitoring glycemic control. Factors affecting accuracy: hemoglobinopathies, anemia, renal failure, pregnancy, and recent transfusions.

\medterm{Hemoglobin Electrophoresis}-- Separation of hemoglobin variants by electrophoresis or HPLC. Normal adult: HbA (95-98\%), HbA2 (2-3\%), HbF (<1\%). Sickle cell disease (HbSS): HbS only, no HbA. Sickle cell trait (HbAS): HbA > HbS. Thalassemia: elevated HbA2 or HbF. Beta thalassemia trait: HbA2 3.5-7\%. HbF elevated in beta thalassemia, sickle cell disease, HPFH, and juvenile myelomonocytic leukemia. Used for diagnosis of hemoglobinopathies.

\medterm{Hepatitis Serology}-- Panel of serologic markers for hepatitis A, B, C, and E. HAV: anti-HAV IgM (acute), anti-HAV IgG (past infection/vaccination). HBV: HBsAg (active infection), anti-HBs (immunity/vaccination), anti-HBc IgM (acute), anti-HBc IgG (past), HBeAg (high replication), anti-HBe (low replication). HBV DNA: quantifies viral load. HCV: anti-HCV (screening, past or present), HCV RNA (active infection, viral load). HEV: anti-HEV IgM/IgG. Used for diagnosis, determination of phase/chronicity, and treatment monitoring.

\medterm{HIV Test}-- Screening and confirmatory testing for HIV. Fourth-generation HIV antigen/antibody combo test (p24 antigen + HIV-1/2 antibodies): window period 2-4 weeks. Positive screening requires confirmatory HIV-1/2 antibody differentiation assay. Rapid point-of-care tests (fingerstick, oral swab): results in 20 minutes. Nucleic acid amplification test (HIV RNA): detects virus as early as 10 days post-infection. Recommended: annual screening for sexually active individuals, all pregnant women, and high-risk populations.

\medterm{Holter Monitor}-- Continuous ECG recording for 24-48 hours to detect arrhythmias, quantify burden (atrial fibrillation, PVCs, PACs), and correlate symptoms with rhythm. Patient diary logs symptoms. Indications: palpitations, syncope, presyncope, and assessment of arrhythmia control. Reveals paroxysmal arrhythmias not captured on standard 12-lead ECG. Extended monitoring (event recorder, loop recorder, mobile cardiac telemetry) for up to 30 days for infrequent symptoms. Implantable loop recorder (Reveal) for annual monitoring.

\medterm{Hydrogen Breath Test}-- Non-invasive test for carbohydrate malabsorption (lactose, fructose) and small intestinal bacterial overgrowth (SIBO). Patient ingests a substrate (lactose, glucose, or lactulose). Breath hydrogen is measured at baseline and q15-30 minutes for 2-3 hours. Rise in H2 >20 ppm from baseline: positive (bacterial fermentation). Lactose intolerance: elevated H2 after lactose challenge, with symptoms. SIBO: early peak (within 60-90 minutes) after glucose or lactulose. False positives: recent antibiotics, colonic fermentation (rapid transit). Antibiotics, colonoscopy prep, and dietary restrictions required before testing.

\medterm{Hysterosalpingography (HSG)}-- Fluoroscopic imaging after injection of iodinated contrast into the uterine cavity through a cervical cannula. Used for infertility evaluation: assesses tubal patency (normal: contrast spills freely into peritoneal cavity), uterine cavity shape (fibroids, polyps, synechiae, Mullerian anomalies), and proximal/distal tubal occlusion. Performed in the follicular phase (days 7-10). Contraindications: pregnancy, active PID, and contrast allergy. Complications: pain, vasovagal reaction, and infection.

\medterm{IGRA (Interferon-Gamma Release Assay)}-- Blood test for Mycobacterium tuberculosis infection (latent or active). T-SPOT.TB (ELISPOT) or QuantiFERON-TB Gold Plus (ELISA). Measures IFN-gamma released by T cells after exposure to MTB-specific antigens (ESAT-6, CFP-10, TB7.7). Advantages over PPD: no cross-reactivity with BCG or most nontuberculous mycobacteria, single patient visit, no booster phenomenon. Preferred for BCG-vaccinated individuals and serial testing. Cannot distinguish latent from active TB.

\medterm{Immunohistochemistry}-- Technique using antibodies to detect specific antigens in tissue sections. Uses: tumor classification and subtyping (Carcinoma: CK7, CK20; Lymphoma: CD markers; Melanoma: S100, HMB45, MelanA; Sarcoma: vimentin, desmin, CD34), determination of tumor origin (CK7+/CK20-: lung, breast; CK7-/CK20+: colorectal), assessment of prognostic markers (ER/PR, HER2 in breast cancer), and identification of infectious agents. Essential in modern pathology for diagnosis and treatment planning.

\medterm{Influenza Test}-- Detection of influenza A and B viruses. Rapid antigen test: nasal/throat swab, results in 15 minutes, specificity >95\%, sensitivity 50-70\% (high false negative rate). RT-PCR: gold standard, highly sensitive and specific, results in 1-4 hours, distinguishes influenza A subtypes (H1N1, H3N2). Used during influenza season for patients with ILI (fever + cough/sore throat). Antiviral therapy (oseltamivir, baloxavir) is most effective when started within 48 hours.

\medterm{Intravenous Pyelogram (IVP)}-- Radiographic imaging of the urinary tract after IV contrast administration. Used to evaluate the kidneys, ureters, and bladder for obstruction, stones, tumors, and anatomic abnormalities. Sequential films show nephrogram (renal parenchyma), pyelogram (collecting system), and ureterogram (ureters). Largely replaced by CT urography and MR urography, which provide better anatomic detail and sensitivity.

\medterm{IV Fluids}-- Intravenous fluid therapy for volume resuscitation and maintenance. Crystalloids: normal saline (0.9\% NaCl, increases chloride, risk of hyperchloremic metabolic acidosis), lactated Ringer (balanced, contains K+, Ca2+, lactate buffer, safer in large volumes), Plasma-Lyte (balanced, closest to plasma). Colloids: albumin (5\%, 25\%), hydroxyethyl starch (HES, nephrotoxic, no longer recommended). Resuscitation: 30 mL/kg bolus for sepsis, hemorrhage. Maintenance: 1.5-2.5 L/day with D5 1/2NS + 20 mEq KCl.

\medterm{Lactate}-- A product of anaerobic metabolism, elevated in tissue hypoperfusion and hypoxia. Normal: <2.0 mmol/L. Lactate >2.0: hyperlactatemia. Lactate >4.0: severe sepsis/septic shock (associated with increased mortality). Types: type A (tissue hypoxia: shock, sepsis, hypoperfusion, seizures, exercise) and type B (no hypoxia: liver disease, thiamine deficiency, metformin, mitochondrial dysfunction, malignancy). Used for sepsis screening (qSOFA), monitoring resuscitation response, and prognostication in critical illness.

\medterm{Lactic Acid}-- See Lactate.

\medterm{Lactose Tolerance Test}-- See Hydrogen Breath Test.

\medterm{Lipid Panel}-- Measurement of total cholesterol, LDL-C, HDL-C, and triglycerides. Fasting (9-12 hours) or non-fasting acceptable for screening. Desirable levels: total cholesterol <200 mg/dL, LDL <100 mg/dL (<70 for very high risk), HDL >40 (men) or >50 (women) mg/dL, triglycerides <150 mg/dL. Used for ASCVD risk assessment (Pooled Cohort Equations), guiding statin therapy, and monitoring treatment. Non-HDL cholesterol (total minus HDL) provides additional risk information.

\medterm{Liver Function Tests}-- A panel of blood tests that assess the health and function of the liver, used to screen for, diagnose, and monitor liver disease and to monitor response to treatment with hepatotoxic drugs. The typical LFT panel includes: (1) Alanine aminotransferase (ALT): a cytosolic enzyme found predominantly in the liver, released into the bloodstream when hepatocytes are damaged. Elevated in: viral hepatitis (all types), alcoholic hepatitis, non-alcoholic fatty liver disease (NAFLD), drug-induced liver injury (DILI), and ischemic hepatitis. ALT is more specific for liver injury than AST (aspartate aminotransferase). AST is present in the liver, heart, muscle, kidney, and brain. The AST/ALT ratio helps distinguish etiology: alcoholic hepatitis (AST:ALT >1.5, usually 2:1), viral hepatitis (AST:ALT <1), and cirrhosis (AST:ALT >1). (2) Alkaline phosphatase (ALP): an enzyme found in bile ducts (canalicular membrane), bone (osteoblasts), placenta, and intestine. Elevated in: cholestatic liver disease (obstructive jaundice, primary biliary cholangitis, primary sclerosing cholangitis), bone disease (Paget, metastasis, fracture, healing), infiltrative disease (sarcoidosis, amyloidosis), and third trimester of pregnancy. Gamma-glutamyl transferase (GGT) is often measured with ALP to distinguish hepatic (GGT elevated) from bone (GGT normal) sources. GGT is induced by alcohol and medications (phenytoin, carbamazepine). (3) Bilirubin: total, conjugated (direct) and unconjugated (indirect). See Jaundice. (4) Albumin: synthesised exclusively by the liver, half-life 20 days. Decreased in chronic liver disease (cirrhosis) indicating impaired synthetic function, nephrotic syndrome, and malnutrition. (5) Prothrombin time (PT) and international normalised ratio (INR): vitamin K-dependent clotting factors (II, VII, IX, X) are synthesised in the liver. PT/INR is prolonged in severe liver disease (hepatic synthetic dysfunction) and vitamin K deficiency. PT/INR is the most sensitive acute measure of liver function in acute liver failure. (6) Total protein: decreased in chronic liver disease, increased in monoclonal gammopathies and chronic inflammation.

\medterm{Liver Function Tests (LFTs)}-- A panel of blood tests assessing liver health: ALT (alanine aminotransferase, primarily hepatocyte injury), AST (aspartate aminotransferase, also released by muscle/heart), ALP (alkaline phosphatase, biliary obstruction, bile duct injury), GGT (gamma-glutamyl transferase, confirms hepatobiliary source of elevated ALP), total and direct bilirubin (jaundice, conjugation/excretion), albumin (synthetic function), and PT/INR (clotting factor synthesis). Patterns: hepatocellular (ALT/AST high), cholestatic (ALP/GGT/bilirubin high), and mixed.

\medterm{Lumbar Puncture}-- See Spinal Tap.

\medterm{Lumbar Puncture (Spinal Tap)}-- Insertion of a needle into the subarachnoid space at L3-L4 or L4-L5 to obtain CSF. Indications: suspected meningitis/encephalitis, subarachnoid hemorrhage, demyelinating disease, and intrathecal therapy. Contraindications: increased intracranial pressure (papilledema, mass lesion on CT), coagulopathy, and infection at the puncture site. Opening pressure: normal 10-20 cm H2O. Complications: post-dural puncture headache (volume depletion, caffeine, blood patch), infection, and bleeding.

\medterm{Mammography}-- X-ray imaging of the breast for cancer screening and diagnosis. Screening mammography: annual for women \(\ge\)40 (USPSTF: biennial 50-74). Views: CC (craniocaudal) and MLO (mediolateral oblique). Digital mammography with CAD (computer-aided detection) improves sensitivity. Tomosynthesis (3D): reduces tissue overlap, improves cancer detection and recall rates. BI-RADS classification: 0 (incomplete), 1 (negative), 2 (benign), 3 (probably benign, short interval), 4 (suspicious, biopsy), 5 (highly suggestive of malignancy), 6 (known biopsy-proven malignancy).

\medterm{Medical Tests}-- Diagnostic procedures used to screen for, confirm, or monitor disease through analysis of biologic specimens, imaging, physiologic measurements, or genetic material. Categories: laboratory tests (complete blood count, comprehensive metabolic panel, cardiac biomarkers, thyroid function, HbA1c, cultures, PCR), imaging studies (X-ray, CT, MRI, ultrasound, PET, nuclear scintigraphy), physiologic testing (ECG, Holter monitor, pulmonary function tests, echocardiography), endoscopic procedures (colonoscopy, esophagogastroduodenoscopy, bronchoscopy, cystoscopy), and genetic/molecular testing (karyotype, microarray, NGS panels). Interpretation requires understanding of sensitivity, specificity, positive/negative predictive values, and reference ranges adjusted for age, sex, and comorbidities. Appropriate use reduces overdiagnosis and incidentalomas; shared decision-making about testing is recommended.

\medterm{Methacholine Challenge Test}-- A bronchoprovocation test for diagnosing asthma when spirometry is normal. Inhaled methacholine (a cholinergic agonist) induces bronchoconstriction. Spirometry is repeated after each escalating dose. Positive: \(\ge\)20\% fall in FEV1 (PC20) at \(\le\)16 mg/mL. Sensitivity >90\% for asthma. Contraindications: FEV1 <50\% predicted, recent MI, uncontrolled hypertension, and pregnancy. Bronchodilator reversibility and asthma treatment affect results.

\medterm{Monospot (EBV Heterophile Antibody)}-- Rapid test for infectious mononucleosis caused by Epstein-Barr virus. Detects heterophile antibodies (IgM) that agglutinate horse or sheep red blood cells. Sensitivity: 70-95\% in symptomatic patients. False negative in early disease (<1 week) and children <5. False positive in autoimmune disease, HIV, lymphoma, and other viral infections. Positive: confirms EBV infection. Negative with high suspicion: order EBV-specific serology (VCA IgM, VCA IgG, EBNA IgG).

\medterm{MRCP (Magnetic Resonance Cholangiopancreatography)}-- Non-invasive MRI of the biliary and pancreatic ducts using heavily T2-weighted sequences that highlight static fluid. Used for: choledocholithiasis (CBD stones), biliary strictures, primary sclerosing cholangitis, pancreatic duct anatomy, and congenital biliary anomalies. Advantages over ERCP: non-invasive, no contrast, no radiation, no pancreatitis risk. Limitations: lower spatial resolution than ERCP, cannot treat (stone removal, stenting). Sensitivity for CBD stones: >90\%.

\medterm{MRI (Magnetic Resonance Imaging)}-- Imaging using strong magnetic fields and radiofrequency pulses to generate detailed cross-sectional images. No ionizing radiation. Superior soft-tissue contrast compared to CT. Sequences: T1-weighted (anatomy, fat bright), T2-weighted (pathology, water bright), FLAIR (suppresses CSF, for brain lesions), DWI (acute stroke, abscess, cellular tumors), and STIR (suppresses fat). Gadolinium contrast: highlights inflammation, infection, and tumors. Contraindications: ferromagnetic implants (some pacemakers, aneurysm clips, metal fragments).

\medterm{MRV (Magnetic Resonance Venography)}-- MRI technique to visualize veins, using time-of-flight (TOF) or contrast-enhanced sequences. Indications: cerebral venous sinus thrombosis, dural sinus stenosis, abdominal/pelvic venous thrombosis, portal vein thrombosis, and renal vein thrombosis. Differentiates acute vs. chronic thrombus. Non-invasive alternative to conventional venography.

\medterm{MUGA Scan (Multigated Acquisition Scan)}-- Radionuclide ventriculography using Tc-99m-labeled RBCs to assess left ventricular function. ECG-gated acquisition synchronized to the QRS complex. Measures: LVEF (most accurate non-invasive method), wall motion, and ventricular volumes. Gold standard for LVEF assessment prior to cardiotoxic chemotherapy (doxorubicin, trastuzumab). Normal LVEF >55\%. Decline >10\% to <50\% indicates cardiotoxicity. Also used for pre-operative risk assessment and heart failure evaluation.

\medterm{Multiplex PCR (BioFire, FilmArray)}-- A rapid molecular diagnostic system using nested multiplex PCR to detect multiple pathogens from a single sample. Panels: respiratory (adenovirus, coronavirus, hMPV, influenza, parainfluenza, RSV, rhinovirus, Bordetella, Chlamydia, Mycoplasma), GI (Campylobacter, Salmonella, Shigella, E. coli, norovirus, rotavirus, etc.), blood culture ID, and meningitis/encephalitis. Results in 45-60 minutes. Sensitivity 90-100\%. Guides early targeted therapy and reduces unnecessary antibiotics.

\medterm{Nerve Conduction Studies}-- Measurement of the speed and amplitude of electrical impulses along peripheral nerves. Sensory NCS: stimulate at one point, record distally. Motor NCS: stimulate at two or more points along the nerve, record from muscle. Parameters: conduction velocity (m/s), amplitude (microvolts/millivolts), and distal latency. Abnormalities: axonal loss (low amplitude, normal velocity), demyelination (slow velocity, prolonged latency, conduction block). Used with EMG to diagnose neuropathy, radiculopathy, and myopathy.

\medterm{Newborn Hearing Screen}-- Universal screening for congenital hearing loss, performed before hospital discharge. Two-step protocol: otoacoustic emissions (OAE: measures cochlear outer hair cell response to click stimulus) and automated auditory brainstem response (AABR: measures neural response from the cochlea to the brainstem). Pass/refer. Refer: requires outpatient diagnostic audiology (full ABR) by 3 months. Early intervention by 6 months improves language outcomes. Risk factors for late-onset/progressive HL: NICU stay >5 days, congenital CMV, family history, and certain syndromes.

\medterm{Newborn Screening}-- A public health program testing newborns for treatable genetic, metabolic, and endocrine disorders 24-48 hours after birth. Blood spot (heel stick) on filter paper (Guthrie card). Conditions vary by state (Recommended Uniform Screening Panel: 35 core conditions). Common: phenylketonuria (PKU), congenital hypothyroidism, cystic fibrosis, sickle cell disease, medium-chain acyl-CoA dehydrogenase (MCAD) deficiency, galactosemia, biotinidase deficiency, congenital adrenal hyperplasia, and hearing loss (by otoacoustic emissions). Positive screens require confirmatory testing and prompt specialist referral. Early treatment prevents morbidity and mortality.

\medterm{Non-Stress Test (NST)}-- Antepartum fetal monitoring assessing fetal heart rate reactivity. Reactive NST: \(\ge\)2 accelerations of 15 bpm for 15 seconds (\(\ge\)32 weeks) or 10 bpm for 10 seconds (<32 weeks) within 20 minutes. Non-reactive: insufficient accelerations after 40 minutes (may indicate fetal sleep, acidosis, or CNS depression). False positive rate high. Used for fetal well-being assessment in high-risk pregnancies (preeclampsia, IUGR, diabetes, postterm, decreased fetal movement). Often combined with AFI (biophysical profile).

\medterm{Pancreatic Elastase}-- Measurement of fecal elastase-1 (FE-1) as a non-invasive test for pancreatic exocrine insufficiency. Values: >200 mcg/g (normal), 100-200 (mild-moderate insufficiency), <100 mcg/g (severe insufficiency). Unaffected by pancreatic enzyme replacement therapy (does not cross-react with porcine enzymes). More sensitive and stable than fecal fat quantification. False positives: watery stool (dilution effect). Used to diagnose chronic pancreatitis, cystic fibrosis, and pancreatic insufficiency.

\medterm{Paracentesis}-- Aspiration of fluid from the peritoneal cavity (ascites). Indications: diagnostic (evaluate new ascites: SAAG, cell count, culture, cytology, albumin, total protein; rule out SBP: PMN >250) and therapeutic (relief of tense ascites, 4-6 L). Site: left lower quadrant (avascular zone), lateral to rectus sheath, ultrasound-guided. Contraindications: overlying cellulitis, severe coagulopathy, and pregnancy. Serum-ascites albumin gradient (SAAG): >1.1 g/dL (portal hypertension: cirrhosis, heart failure), <1.1 g/dL (peritoneal disease: malignancy, TB, pancreatitis, nephrotic). Large-volume paracentesis requires albumin infusion (6-8 g/L removed) to prevent post-paracentesis circulatory dysfunction.

\medterm{PEG Tube (Percutaneous Endoscopic Gastrostomy)}-- Endoscopic placement of a feeding tube through the abdominal wall directly into the stomach. Indications: long-term enteral feeding (>4 weeks) for dysphagia (stroke, ALS, dementia, head/neck cancer) and malnutrition. Contraindications: gastric outlet obstruction, severe ascites, coagulopathy, gastric cancer, and inability to transilluminate the abdominal wall. Procedure: EGD with transillumination, needle puncture under direct endoscopic visualization, wire passage, and tube placement (Pull or Push technique). Replacements: 8-12 months later (balloon or button type). Complications: aspiration, infection (peristomal cellulitis), bleeding, perforation, and tube dislodgement. Buried bumper syndrome: overgrowth of gastric mucosa over the internal bumper.

\medterm{Pericardiocentesis}-- Aspiration of pericardial fluid using a needle inserted through the subxiphoid or apical approach, under echocardiographic or fluoroscopic guidance. Indications: cardiac tamponade (therapeutic emergency), diagnostic evaluation of pericardial effusion (malignancy, infection, autoimmune). Complications: myocardial puncture/laceration, coronary artery injury, pneumothorax, arrhythmia, and infection. A small catheter can be left in place for continuous drainage. Fluid analysis: cell count, cytology, culture, TB PCR/ culture, and rheumatologic studies. High-risk procedure; should be performed with surgical backup.

\medterm{Periodontal Probing}-- Measurement of the depth of the gingival sulcus/pocket using a periodontal probe with millimeter markings. Normal: 1-3 mm. Periodontal pockets: \(\ge\)4 mm (indicates attachment loss and periodontitis). Full-mouth probing at 6 sites per tooth: mesio-buccal, mid-buccal, disto-buccal, mesio-lingual, mid-lingual, disto-lingual. Bleeding on probing (BOP): sign of active inflammation. Clinical attachment level (CAL): distance from CEJ to base of pocket. Used for periodontal disease diagnosis, staging (Stage I-IV), and monitoring treatment response.

\medterm{Peripheral Blood Smear}-- Microscopic examination of a stained blood smear (Wright-Giemsa stain). Evaluates: RBC morphology (size: macrocytic, microcytic; shape: spherocytes, sickle cells, schistocytes, target cells, elliptocytes; color: hypochromic, polychromasia; inclusions: basophilic stippling, Howell-Jolly bodies, Pappenheimer bodies), WBC morphology (blast cells, band forms, toxic granulation, Doehle bodies, smudge cells, Auer rods), and platelet morphology (size, granules). Performed when automated CBC shows abnormalities.

\medterm{PET Scan (Positron Emission Tomography)}-- Nuclear medicine imaging using radioactive tracers (most commonly FDG, a glucose analog) to detect metabolically active tissues. Indications: cancer staging and restaging (lung, lymphoma, melanoma, breast, colon), evaluation of solitary pulmonary nodule, myocardial viability (FDG uptake in hibernating myocardium), and neurologic (Alzheimer disease: hypometabolism; epilepsy: seizure focus). Standardized uptake value (SUV) quantifies tracer activity. Combined PET/CT provides anatomic localization. False positives in infection/inflammation.

\medterm{Polysomnography (Sleep Study)}-- Overnight monitoring for sleep-disordered breathing. Parameters: EEG (sleep staging: N1, N2, N3, REM), EOG (eye movements), EMG (chin and leg movements), airflow (nasal pressure, thermistor), respiratory effort (chest/abdomen belts), snore microphone, body position, oximetry (SpO2), and ECG. Calculates AHI (apnea-hypopnea index): normal <5, mild 5-15, moderate 15-30, severe >30. Used to diagnose obstructive sleep apnea, central sleep apnea, periodic limb movement disorder, and narcolepsy (MSLT following PSG).

\medterm{Pregnancy Test}-- See hCG (Human Chorionic Gonadotropin).

\medterm{Preoperative Evaluation}-- Medical assessment before elective surgery to optimize outcomes and minimize perioperative risk. Components: history and physical (comorbidities: CAD, CHF, COPD, diabetes, CKD), medications (anticoagulants, antiplatelets, antihyperglycemics), functional status (METs), and cardiac risk assessment (RCRI: Revised Cardiac Risk Index, NSQIP MICA risk calculator). Testing per RCRI and age/symptoms: ECG (age >65, cardiac risk factors, known CAD), CBC, BMP, coagulation studies, type and screen (for planned transfusions), CXR (age >70, cardiac/pulmonary disease), and PFTs (lung resection, COPD exacerbation). Nothing routine; testing guided by specific risk factors and surgical severity.

\medterm{Procalcitonin}-- A biomarker for bacterial infection and sepsis, produced by thyroid C-cells and peripheral tissues in response to bacterial toxins. Levels: <0.1 ng/mL (normal), 0.1-0.25 (unlikely bacterial), 0.25-0.5 (possible), >0.5 (probable bacterial infection), >2.0 (severe sepsis/septic shock). More specific for bacterial infection than CRP. Used to guide antibiotic initiation and discontinuation (procalcitonin algorithms reduce antibiotic duration). Elevated in bacterial pneumonia, meningitis, UTI. Not elevated in viral infections or non-infectious inflammation.

\medterm{PSA (Prostate-Specific Antigen)}-- A serine protease produced by the prostate, used for prostate cancer screening. Normal: <4.0 ng/mL, but age-specific ranges exist. PSA >10 ng/mL: high probability of cancer. PSA 4-10: gray zone; percent free PSA (<25\% suggests cancer) and PSA density help. Screening controversy: reduces prostate cancer mortality (1.3 per 1000) but leads to overdiagnosis/overtreatment. Shared decision-making for men 55-69. Annual screening recommended for high-risk (African American, family history).

\medterm{Pulmonary Function Tests}-- Spirometry: FVC (forced vital capacity), FEV1 (forced expiratory volume in 1 second), FEV1/FVC ratio. Obstructive pattern: FEV1/FVC <0.7 (asthma, COPD). Restrictive pattern: FVC <80\% with normal or increased FEV1/FVC (pulmonary fibrosis, chest wall disorders). Lung volumes: TLC (total lung capacity), RV (residual volume). Diffusion capacity (DLCO): measures gas transfer across the alveolar-capillary membrane. Indications: dyspnea evaluation, COPD staging (GOLD), pre-operative assessment, and monitoring lung disease progression.

\medterm{Quad Screen (Maternal Serum Screening)}-- Second-trimester prenatal screening (15-22 weeks) for fetal aneuploidy and neural tube defects. Measures 4 markers: AFP (alpha-fetoprotein: elevated in NTD, low in Down syndrome), hCG (elevated in Down syndrome), uE3 (unconjugated estriol: low in Down syndrome), and inhibin A (elevated in Down syndrome). Detection rates: trisomy 21 (80\%), trisomy 18 (60-70\%), NTD (80\%). Replaced in many settings by cell-free DNA (cfDNA, NIPT) with higher detection rates.

\medterm{Rapid Strep Test}-- Immunoassay for Group A Streptococcus (Streptococcus pyogenes) from throat swab. Results: 5-10 minutes. Specificity: >95\%. Sensitivity: 80-90\% compared to culture. Positive: treat. Negative: confirm with throat culture (especially in children). Used for diagnosis of GAS pharyngitis to guide antibiotic therapy (prevents unnecessary antibiotics for viral pharyngitis). First-line treatment for confirmed GAS: penicillin or amoxicillin.

\medterm{Reticulocyte Count}-- Percentage of young RBCs in the peripheral blood, reflecting bone marrow erythroid activity. Normal: 0.5-2.5\% (absolute 25-100 K/uL). Reticulocyte production index (RPI) corrects for anemia severity: RPI = reticulocyte \% x (Hct/45) / maturation correction factor. High RPI (>2-3): bone marrow responding appropriately (hemolysis, acute blood loss). Low RPI (<2): bone marrow not responding (iron deficiency, aplastic anemia, marrow infiltration). Used to classify anemia as hyperproliferative vs. hypoproliferative.

\medterm{Rheumatoid Factor (RF)}-- An autoantibody directed against the Fc portion of IgG, positive in 70-80\% of rheumatoid arthritis patients. Also positive in other autoimmune diseases (SLE, Sjogren, mixed cryoglobulinemia), chronic infections (TB, endocarditis, hepatitis C), and 5\% of healthy elderly. Low sensitivity (may be negative in seronegative RA) and moderate specificity. Used in ACR/EULAR classification criteria for RA with anti-CCP. Serial monitoring is not useful. Higher titers associated with more severe, erosive disease.

\medterm{RPR/VDRL (Syphilis Serology)}-- Non-treponemal screening tests for syphilis. Nontreponemal: RPR (rapid plasma reagin) or VDRL (venereal disease research laboratory). Positive in primary syphilis (1-4 weeks after chancre), secondary (100\% positive), and early latent. Quantitative: titer correlates with disease activity and treatment response (fourfold decline indicates successful treatment). False positives: pregnancy, SLE, TB, IV drug use, hepatitis, HIV. Positive screening requires confirmatory treponemal test (FTA-ABS, TP-PA, EIA). Titer monitoring after treatment.

\medterm{Serum Protein Electrophoresis (SPEP)}-- Separation of serum proteins by electrophoresis into 5 fractions: albumin, alpha-1, alpha-2, beta, and gamma. Monoclonal gammopathy: M-spike in gamma region (multiple myeloma, MGUS, Waldenstrom macroglobulinemia). Immunofixation: identifies heavy chain (IgG, IgA, IgM, IgD, IgE) and light chain (kappa, lambda) isotype. Serum free light chain assay: kappa/lambda ratio (abnormal in light chain myeloma, AL amyloidosis). Used to screen for, diagnose, and monitor plasma cell dyscrasias.

\medterm{SPECT Scan (Single Photon Emission CT)}-- Nuclear medicine imaging using gamma-emitting radioisotopes (Tc-99m, Tl-201, I-123) with a rotating gamma camera to create 3D images. Indications: myocardial perfusion imaging (MPI with sestamibi, dobutamine or regadenosone stress), bone scan (metastases, stress fracture, osteomyelitis), brain perfusion (dementia, seizure focus), and VQ scan (Tc-99m MAA for perfusion, Xe-133 or Tc-DTPA for ventilation). Combined with CT (SPECT/CT) improves anatomic localization.

\medterm{Sweat Chloride Test}-- The gold standard diagnostic test for cystic fibrosis. Measures chloride concentration in sweat after pilocarpine iontophoresis. Normal: <30 mmol/L. Intermediate: 30-59. Positive: \(\ge\)60 mmol/L. False positives: malnutrition, adrenal insufficiency, hypothyroidism, and ectodermal dysplasia. False negatives: rare CFTR mutations with normal sweat chloride. Requires adequate sweat collection (>75 mg). Part of CF diagnostic algorithm with clinical features and CFTR genetic testing.

\medterm{Thyroid Function Tests}-- Measurement of TSH (most sensitive screening test), free T4, and free T3. TSH high, T4 low $\rightarrow$ primary hypothyroidism. TSH low, T4 high $\rightarrow$ primary hyperthyroidism. TSH low, T4 normal $\rightarrow$ subclinical hyperthyroidism or T3-toxicosis (measure T3). TSH high, T4 normal $\rightarrow$ subclinical hypothyroidism. TSH normal, T4 low $\rightarrow$ central hypothyroidism (pituitary/hypothalamic). Anti-TPO antibodies: positive in Hashimoto thyroiditis. TSH receptor antibodies (TRAb): positive in Graves disease.

\medterm{Tilt Table Test}-- Evaluation of unexplained syncope by reproducing orthostatic stress. Patient is strapped to a motorized table tilted from supine to 60-80 degrees for 20-45 minutes. Positive: reproduction of symptoms with hypotension and/or bradycardia (vasovagal syncope, delayed orthostatic hypotension). Also assesses for postural orthostatic tachycardia syndrome (POTS: HR increase >30 bpm without hypotension). Isoproterenol or nitroglycerin provocation increases sensitivity. Used when initial evaluation for syncope is inconclusive.

\medterm{TORCH Panel}-- Serologic screening for congenital infections: Toxoplasma, Rubella, CMV, Herpes simplex (and others: syphilis, parvovirus B19, VZV, Zika). IgM: indicates recent/active infection (maternal or fetal). IgG: past infection or passive immunity. Paired maternal and neonatal testing. Indications: pregnancy with febrile illness or rash, fetal abnormalities on ultrasound (hydrops, microcephaly, calcifications), and neonatal evaluation for congenital infection. Management depends on the specific pathogen identified.

\medterm{Troponin}-- Cardiac-specific troponin I and T, the preferred biomarkers for myocardial injury (MI). Highly sensitive and specific. Rise 2-4 hours after symptom onset, peak at 12-24 hours, remain elevated for 7-10 days. High-sensitivity troponin (hs-cTn) allows earlier detection (1-2 hours). Fourth Universal Definition: detection of a rise/fall of troponin with at least one value above the 99th percentile upper reference limit, plus clinical evidence of myocardial ischemia. Used for diagnosis of NSTEMI, risk stratification in ACS, and prognosis in many cardiac conditions.

\medterm{Tuberculin Skin Test (PPD)}-- Intradermal injection of 0.1 mL of purified protein derivative (PPD) on the forearm (Mantoux method). Read at 48-72 hours measuring induration (not erythema). Positive thresholds: \(\ge\)5 mm (HIV, close contacts, transplant, immunosuppressed), \(\ge\)10 mm (recent immigrants, IV drug users, healthcare workers, children <4, high-risk settings), \(\ge\)15 mm (no risk factors). Does not differentiate latent from active TB. False positives: BCG vaccination, nontuberculous mycobacteria. Positive test requires chest X-ray and symptom evaluation.

\medterm{Ultrasonography}-- See Ultrasound (Ultrasonography).

\medterm{Ultrasound}-- Imaging using high-frequency sound waves (2-18 MHz). No ionizing radiation. Right upper quadrant (RUQ): gallbladder (gallstones, cholecystitis), liver (steatosis, cysts), kidneys, and pancreas. Pelvic: uterus, ovaries, adnexa. Renal: hydronephrosis, stones, cysts, tumors. Scrotal: testicular torsion, epididymitis, varicocele, hydrocele. Thyroid: nodules (TI-RADS classification), goiter. Soft tissue: abscess, cellulitis, hematoma, lipoma. Vascular: DVT, carotid stenosis, AAA. Point-of-care (POCUS): FAST, cardiac, lung.

\medterm{Umbilical Cord Blood Gas}-- Analysis of arterial and venous cord blood collected immediately after delivery. Umbilical artery (UA): reflects fetal (baby's) acid-base status. Normal UA pH >7.20, base excess < $-$10. Umbilical vein (UV): reflects maternal/placental status. Indications: suspected fetal distress, low Apgar scores, and medicolegal documentation of birth asphyxia. Fetal acidosis: pH <7.0 (severe) or pH 7.0-7.19 (moderate). Mixed acidosis (respiratory + metabolic component) indicates prolonged hypoxic event.

\medterm{Upper Endoscopy}-- See Esophagogastroduodenoscopy (EGD).

\medterm{Upper GI Series}-- Fluoroscopic examination of the esophagus, stomach, and duodenum after oral barium contrast. Single-contrast: looks for gross abnormalities (obstruction, masses). Double-contrast (barium + gas): better mucosal detail for ulcers, erosions, and early malignancies. Indications: dysphagia, dyspepsia, nausea/vomiting, weight loss, and GI bleeding. Can assess for hiatal hernia, ulcers, gastric cancer, and outlet obstruction. Largely replaced by EGD.

\medterm{Urinalysis}-- Analysis of urine. Dipstick: specific gravity, pH, leukocyte esterase (WBCs), nitrite (bacteria), protein, glucose, ketones, bilirubin, urobilinogen, blood, and hemoglobin. Microscopy: RBCs, WBCs, casts (hyaline, granular, RBC, WBC), crystals (calcium oxalate, uric acid, triple phosphate), bacteria, yeast, squamous cells, and epithelial cells. Indications: UTI, hematuria, renal disease (proteinuria, casts), and screening. Clean-catch midstream specimen for culture reduces contamination.

\medterm{Urine Drug Screen}-- Immunoassay-based screening for common drugs of abuse in urine. Detects: amphetamines, benzodiazepines, cocaine (benzoylecgonine), opiates (morphine, codeine, heroin), oxycodone, methadone, fentanyl, barbiturates, PCP, THC (cannabinoids), and alcohol. Pros: rapid (minutes at POC), inexpensive. Cons: false positives (poppy seeds $\rightarrow$ opiates, certain antihistamines $\rightarrow$ PCP, NSAIDs $\rightarrow$ barbiturates), false negatives (synthetic cannabinoids/spice, synthetic cathinones/bath salts, LSD). Positive screen requires confirmatory GC-MS or LC-MS/MS.

\medterm{Urine hCG (Pregnancy Test)}-- Rapid point-of-care test detecting human chorionic gonadotropin in urine. Sensitivity: detects as low as 25 mIU/mL. Positive as early as 12-14 days after conception (before missed period). Results in 3-5 minutes. False negative: dilute urine, very early pregnancy, ectopic pregnancy. False positive: recent pregnancy, hCG-producing tumors (GTN, germ cell), or rarely due to cross-reacting substances. Most sensitive and specific when performed on first morning urine specimen.

\medterm{Venesection (Phlebotomy)}-- The controlled removal of blood from a vein for therapeutic or diagnostic purposes. Therapeutic phlebotomy: treatment for haemochromatosis (removal of excess iron by reducing red cell mass and stimulating erythropoiesis, which mobilises iron stores), polycythaemia vera, and porphyria cutanea tarda; typical removal of 450-500 mL per session, frequency determined by haematocrit or ferritin levels. Diagnostic phlebotomy: blood collection (venepuncture) for laboratory testing, using evacuated tube system (Vacutainer) with appropriate anticoagulant or additive tubes (EDTA for hematology, citrate for coagulation studies, SST/gel for chemistry, fluoride oxalate for glucose). Site: median cubital vein (antecubital fossa) most common. Complications: haematoma, nerve injury (rare), infection, vasovagal syncope, and excessive bleeding in coagulopathic patients.

\medterm{Venipuncture}-- Puncture of a vein to obtain blood specimens or administer IV therapy. Common sites: median cubital vein (most common), cephalic vein, basilic vein, and dorsal hand veins. Steps: tourniquet application (30-60 seconds), vein palpation, skin antisepsis (alcohol, chlorhexidine), needle insertion at 15-30 degree angle, blood collection into appropriate tubes (order of draw: culture first, then serum, then citrate, then heparin, then EDTA, then tubes with additives), and proper needle disposal. Complications: hematoma, hemolysis, nerve injury, infection, and needlestick injury.

\medterm{Viral Load (HIV RNA)}-- Quantitative measurement of HIV-1 RNA in plasma by RT-PCR (Roche COBAS, Abbott RealTime). Normal: undetectable (<20-40 copies/mL). Treatment goal: undetectable viral load (U=U: Undetectable = Untransmittable). Initiation: start ART regardless of viral load or CD4. Monitoring: at baseline, 4-6 weeks after ART start, then q3-6 months. Virologic failure: persistently detectable or rebound >200 copies/mL. Resistance testing before ART switch.

\medterm{Vitamin B12}-- Cobalamin, essential for DNA synthesis, red blood cell production, and neurologic function. Normal: 200-900 pg/mL. Deficiency: <200 pg/mL (but borderline 200-300 may still be symptomatic). Causes: pernicious anemia (autoimmune, most common), gastrectomy, Crohn disease, ileal resection, strict vegan diet, proton pump inhibitors, and metformin. Symptoms: megaloblastic anemia, neuropathy (paresthesias, ataxia), cognitive decline, and glossitis. Methylmalonic acid and homocysteine are more sensitive markers for B12 deficiency.

\medterm{Vitamin D}-- 25-hydroxyvitamin D, the storage form of vitamin D. Normal: >30 ng/mL (75 nmol/L). Insufficiency: 20-30 ng/mL. Deficiency: <20 ng/mL (<50 nmol/L). Extremely deficient: <10 ng/mL (high risk of osteomalacia/rickets). Sources: sun exposure (80\%), diet, and supplements. Deficiency causes secondary hyperparathyroidism, increased bone turnover, osteoporosis, and increased fall/fracture risk. Associated with autoimmune disease, cardiovascular disease, and cancer (causal role unproven). Supplementation: 800-2000 IU/day for deficiency.

\medterm{Voiding Cystourethrogram (VCUG)}-- Fluoroscopic imaging of the bladder and urethra during voiding, after filling the bladder with contrast via a catheter. Indications: evaluation of vesicoureteral reflux (VUR) in children with recurrent UTIs, posterior urethral valves in males, and post-surgical assessment of the bladder/urethra. Grading of VUR: I (ureter only), II (ureter and pelvis, no dilation), III (mild dilation), IV (moderate dilation, tortuous ureter), V (severe dilation, tortuous ureter, intrarenal reflux). Also evaluates bladder capacity, wall thickness, and post-void residual.

\medterm{X-Ray (Plain Radiography)}-- The most common form of medical imaging, using X-rays to create a 2D projection image. Uses: chest (pneumonia, pneumothorax, effusion, cardiomegaly, mass), abdomen (obstruction, perforation, stones), musculoskeletal (fractures, dislocation, arthritis, osteomyelitis), and spine (fracture, alignment, degeneration). Views: PA (posteroanterior) and lateral chest, AP and lateral for spine, AP and frog-leg for hips. Density: air (black), fat (dark gray), soft tissue/water (gray), calcium/bone (white), metal (bright white).

\medterm{X-Rays (Röntgen Rays)}-- See X-Ray (Plain Radiography).
